\section{Preguntas del laboratorio}
\subsection{Parte A}
\subsubsection{Circuito combinacional, secuencial y proceso síncrono}
La distinción principal fue la presencia de memoria: un circuito combinacional entregó una función de las entradas presentes y uno secuencial conservó un estado entre ciclos. En VHDL, \texttt{rising\_edge(clk)} delimitó el instante de actualización de los registros. Con reset síncrono, su acción quedó sujeta al mismo flanco, lo que hizo posible un comportamiento temporal determinista.

\subsubsection{MAX para 1 Hz y 1 kHz y anchura del contador}
Aplicando la Ec.~\eqref{eq:max}, se obtuvo \(\mathrm{MAX}_{1\,\mathrm{Hz}}=100\,000\,000\) y \(\mathrm{MAX}_{1\,\mathrm{kHz}}=100\,000\). Para representar los valores comprendidos entre cero y \(99\,999\,999\), fueron suficientes 27 bits porque \(2^{26}=67\,108\,864<100\,000\,000\leq 2^{27}=134\,217\,728\).

\subsubsection{Habilitación en lugar de reloj dividido}
El empleo de \texttt{clock enable} mantuvo el contador y el divisor sincronizados con el mismo reloj físico. Así se evitó tratar una señal generada por lógica ordinaria como una nueva red de reloj, lo que habría complicado la distribución, las restricciones temporales y el análisis de cruces entre dominios. La operación de un registro quedó condicionada a \texttt{en='1'} sin alterar su fuente de reloj.

\subsubsection{Uso de genéricos en simulación}
Con el valor real de cien millones de ciclos por pulso, una simulación de varios incrementos implicaría una cantidad muy elevada de eventos, tal como se explicó en el caso de estudio del material de divisores \cite{cursodiv}. El parámetro genérico \texttt{MAX} permitió sustituir el valor de implementación por uno pequeño, como cuatro, sin modificar la estructura lógica. En el testbench de A fueron utilizados MAX=4 y N=3, y las comprobaciones finalizaron a los 455,001 ns.

\subsubsection{Frecuencia de conmutación de los LEDs}
Para incrementos a 1 Hz, el bit \(i\) del contador cambió de nivel cada \(2^i\) segundos, y completó un período de encendido y apagado cada \(2^{i+1}\) segundos. En consecuencia, su frecuencia de señal cuadrada fue
\begin{equation}
f_{\mathrm{LED}i}=\frac{1\,\mathrm{Hz}}{2^{i+1}}.
\end{equation}
Por ejemplo, \texttt{LED0} tuvo una frecuencia de señal de 0,5 Hz (una transición por segundo), \texttt{LED1} de 0,25 Hz y \texttt{LED9} de \(1/1024\) Hz. Esta distinción evitó confundir la frecuencia de incremento del contador con la frecuencia de parpadeo de cada LED.

\subsection{Parte B}
\subsubsection{Definición y componentes de la FSM}
La FSM se representó mediante un registro de estado, una lógica de transición y una función de salida. Los cuatro estados de la Fig.~\ref{fig:fsm} representaron, en orden, AND, OR, XOR y ADD. Cada habilitación especificó el paso siguiente; después de ADD se regresó a AND. El reset estableció el estado inicial AND.

\subsubsection{Diferencia entre Moore y Mealy}
En una máquina de Moore, la salida codificada depende exclusivamente del estado almacenado; en una máquina de Mealy, también puede depender de las entradas actuales. La FSM propuesta fue de Moore porque \texttt{op} se obtuvo del estado. Aunque los resultados del datapath variaron combinacionalmente al modificarse A o B, ello no convirtió a la FSM de control en una máquina de Mealy.

\subsubsection{Estados enumerados y síntesis}
La descripción mediante un tipo enumerado permitió utilizar nombres semánticos y redujo errores asociados a códigos binarios asignados manualmente. En general, la síntesis transforma los estados simbólicos a una representación física apropiada y puede escoger una codificación distinta de la codificación de salida \texttt{op}, según sus opciones de optimización. Por tanto, la numeración de \texttt{op} no implicó necesariamente que el registro físico de estado utilizara exactamente esos dos bits.

\subsubsection{Separación de procesos síncronos y combinacionales}
El registro de estado se planteó en un proceso síncrono para conservar el estado entre flancos y actualizarlo solamente al recibirse una habilitación. La lógica de próximo estado se planteó combinacionalmente para evaluar las reglas de transición en función del estado actual. La salida de Moore se definió como función combinacional de ese estado. En el registro de la FSM fue utilizada la condición \texttt{elsif en='1' then estado <= siguiente;} dentro de \texttt{rising\_edge(clk)}, con prioridad del reset. En procesos \texttt{process(all)} separados fueron definidas las transiciones mediante \texttt{case estado} y la decodificación de \texttt{op}. De ese modo no se crearon registros adicionales en la lógica combinacional.

\subsubsection{Relación con la arquitectura de un procesador}
La FSM cumplió una función de control al seleccionar la operación mediante \texttt{op}. El datapath realizó las transformaciones de datos sobre A y B y generó los resultados. Esta separación fue análoga, en una escala sencilla, a la organización de un procesador: la unidad de control determina la acción y la ruta de datos la ejecuta. La arquitectura del laboratorio no constituyó un procesador completo, ya que no incorporó un programa de instrucciones ni todos sus registros y rutas de acceso.

